% Editable manuscript methods, matched to docs/methods and the released stack.
% Requires amsmath and amssymb. Include with % Editable manuscript methods, matched to docs/methods and the released stack.
% Requires amsmath and amssymb. Include with % Editable manuscript methods, matched to docs/methods and the released stack.
% Requires amsmath and amssymb. Include with % Editable manuscript methods, matched to docs/methods and the released stack.
% Requires amsmath and amssymb. Include with \input{methods.tex}.
% Implementation references and full typed-energy parameter tables are in docs/methods.
\section{Methods}
\label{sec:effdock-methods}

\subsection{Task, data and information boundary}
EFF-Dock predicts ligand poses given a receptor, ligand chemistry and a
supplied pocket centre. The experiments are retrospective supplied-pocket
redocking: pocket centres may be reference-ligand-derived and receptors may
be holo structures. The reference pose supplies supervision, mapping and
evaluation labels; its coordinates do not enter the learned vector field,
confidence ranker or minimized refinement objective.

Training structures originate from PLINDER 2024-06/v2. A sample is a system
and ligand-instance pair. The preserved split contains 47,310 training and
1,076 validation IDs. The released docking fine-tune loader contains 47,277
samples; confidence uses 43,092 training and 1,035 validation samples. Their
training intersection is 43,067, with 4,210 docking-only and 25 confidence-only
samples. These independently filtered inventories are not nested and do not
represent the union of every predecessor checkpoint's training exposure.
Public membership records provide IDs, exclusions and SHA-256 hashes.

The historical external exclusion combined a ligand-match condition with a
pocket-community-match condition whose reference communities were identified
from benchmark PDB hits in the processed pool. It is distinct from the newer
strict split builder. No post-hoc exclusion changes the released benchmarks.

\subsection{Fragment representation and equivariant model}
Ligands are partitioned at eligible single, non-ring, non-planar-conjugated
rotatable bonds, retaining at least two heavy atoms per fragment after
singleton merging. Let $f(a)$ assign atom $a$ to fragment $f$, $T_f$ be its
centroid, $q_f$ a unit quaternion and $\ell_a$ its fixed local coordinate:
\begin{equation}
 T_f=\frac{1}{|f|}\sum_{a\in f}x_a,\qquad
 x_a(T,q)=R(q_{f(a)})\ell_a+T_{f(a)}.
 \label{eq:effdock-reconstruction}
\end{equation}
The heterogeneous graph contains ligand atoms, fragments, protein atoms and
residue virtual nodes. Static topology and dynamic protein--ligand contacts
within $5\,\text{\AA}$ condition six equivariant interaction layers, using
32 radial basis functions and spherical harmonics through degree two. Hidden
irreps are $384\times0e+32\times1o+32\times1e+16\times2e+16\times2o$,
with flattened width 736. Irrep notation alone does not establish full
reflection equivariance of stereochemical inputs; the pose contract uses
proper rotations and joint SE(3) transformations.

The model outputs atom vector fields $g_a$. Define $r_a=x_a-T_{f(a)}$.
Unit-weight Newton--Euler aggregation gives
\begin{align}
 v_f&=\frac{1}{|f|}\sum_{a\in f}g_a,&
 \tau_f&=\sum_{a\in f}r_a\times g_a,\\
 I_f&=\sum_{a\in f}(\|r_a\|^2I_3-r_ar_a^{\mathsf T}),&
 \omega_f&=I_f^+\tau_f.
 \label{eq:effdock-head}
\end{align}
The pseudoinverse retains inertia eigenvalues strictly above 1\% of the
largest eigenvalue, yielding an observable-rotation projector $P_f$.

\subsection{Conditional flow and training objective}
For supplied centre $c$, translation priors are
$T_f^0=c+\sigma\epsilon_f$, $\epsilon_f\sim\mathcal N(0,I_3)$;
initial orientations are uniform on SO(3). Conditional interpolation is
\begin{align}
 T_f^t&=(1-t)T_f^0+tT_f^1,&
 q_f^t&=\operatorname{SLERP}(q_f^0,q_f^1;t),\\
 v_f^*&=T_f^1-T_f^0,&
 \omega_f^*&=\operatorname{axisAngle}(q_f^1\otimes(q_f^0)^{-1}).
 \label{eq:effdock-flow-targets}
\end{align}
Quaternion updates use left multiplication. With $F$ fragments and
$F_\omega$ fragments having more than one atom, the flow terms are
\begin{equation}
 L_v=\frac{1}{3F}\sum_f\|v_f-v_f^*\|^2,\qquad
 L_\omega=\frac{1}{3F_\omega}\sum_{f:|f|>1}
 \|\omega_f-P_f\omega_f^*\|^2.
 \label{eq:effdock-flow-loss}
\end{equation}
A zero eligible count yields zero angular loss. The angular denominator is
not the number of observable eigenvectors. The released fragmenter avoids
singletons. Optional time weighting of these flow terms is disabled.

The atom auxiliary loss compares instantaneous velocities, not endpoint
coordinates. For $r_a^t=x_a^t-T_{f(a)}^t$,
\begin{equation}
 u_a=v_{f(a)}+\omega_{f(a)}\times r_a^t,\quad
 u_a^*=v_{f(a)}^*+\omega_{f(a)}^*\times r_a^t,\quad
 L_{\rm atom}=\frac{1}{3A}\sum_a\|u_a-u_a^*\|^2.
 \label{eq:effdock-atom-loss}
\end{equation}
The auxiliary target uses the unprojected angular field; an exactly
unobservable rotation induces zero atom motion.

A separate distance-geometry term uses the one-step endpoint
\begin{equation}
 \widehat T_f^1=T_f^t+(1-t)v_f,\qquad
 \widehat q_f^1=\operatorname{Exp}_q((1-t)\omega_f)\otimes q_f^t.
\end{equation}
Here $\operatorname{Exp}_q$ converts an axis-angle vector to a quaternion.
For unordered inter-fragment atom pairs $\mathcal P_b$ in complex $b$,
\begin{align}
 D_b&=\frac{1}{|\mathcal P_b|}\sum_{(a,c)\in\mathcal P_b}
 \left(\|\widehat x_a^1-\widehat x_c^1\|-\|x_a^1-x_c^1\|\right)^2,\\
 L_{\rm DG}&=\frac{\sum_{b:|\mathcal P_b|>0}t_b^2D_b}
 {\sum_{b:|\mathcal P_b|>0}t_b^2},\qquad
 L=L_v+8L_\omega+0.3L_{\rm atom}+3L_{\rm DG}.
 \label{eq:effdock-total-loss}
\end{align}
A zero denominator yields zero $L_{\rm DG}$. Distributed reductions reproduce
global fragment, atom and weighted-complex means for the respective terms;
these are not equal-weight-per-complex reductions throughout.

The released 50k fine-tune samples
\begin{align}
 p(t)&=0.80p_{\rm SF}(t)+0.10U(0,0.3)+0.10\delta_0,\\
 p_{\rm SF}&=0.02U(0,1)+0.98\operatorname{Law}
 [\operatorname{sigmoid}(0.8+1.7Z)],\quad Z\sim\mathcal N(0,1).
\end{align}
It uses fresh AdamW state, global batch 64 on four GPUs, peak learning rate
$2\times10^{-5}$, weight decay 0.01, gradient clipping 1, and EMA decay 0.999.
Warmup lasts 1,000 updates; cosine decay begins at 40,000 and reaches
$2\times10^{-6}$ at 50,000. The released docking state is terminal EMA.

\subsection{Sampling and confidence}
Main sampling uses $N=100$ candidates, $S=10$ steps, $\sigma=2\,\text{\AA}$
and a $10\,\text{\AA}$ crop. The late-power-three Euler grid and updates are
\begin{align}
 t_j&=1-(1-j/S)^3,\quad j=0,\ldots,S,\\
 T_f^{j+1}&=T_f^j+\Delta t_jv_{\theta,f}(t_j),\qquad
 q_f^{j+1}=\operatorname{Exp}_q(\Delta t_j\omega_{\theta,f}(t_j))\otimes q_f^j.
 \label{eq:effdock-sampler}
\end{align}
Raw and refined banks preserve generation indices and are independently scored.

The confidence scorer uses four equivariant layers, paired docking features,
and attention pooling over global and contact features. It predicts log-one-plus pose
RMSD $\widehat y_k$, pose-success logit $z_k$, atom log displacement
$\widehat y_{ka}$ and atom-success logit $z_{ka}$. Atom and pose regression
use Huber loss with transition one; classification uses BCE with logits and
strict $2\,\text{\AA}$ labels. Let $y_k=\log(1+r_k)$:
\begin{align}
 L_{\rm rank}&=\underset{y_j-y_i>0.3}{\operatorname{mean}}
 \max(0,0.05-(\widehat y_j-\widehat y_i)),\\
 q_k&=\operatorname{sigmoid}((2-r_k)/0.2),\quad
 p_k=\frac{q_k}{\sum_jq_j},\qquad
 L_{\rm list}=-\sum_kp_k\log\operatorname{softmax}(z)_k,\\
 L_{\rm conf}&=0.2L_{\rm atom}+0.2L_{\rm atom\mbox{-}ok}
 +0.3L_{\rm pose}+0.4L_{\rm pose\mbox{-}ok}
 +0.1L_{\rm rank}+L_{\rm list}.
 \label{eq:effdock-confidence}
\end{align}
Here the confidence atom term is displacement regression, distinct from the
docking atom-velocity term in Eq.~\ref{eq:effdock-atom-loss}. Empty eligible
pairs give zero rank loss. Numerically zero quality sums use a uniform target;
one-candidate pools have zero listwise loss.

Training banks contain 32 raw poses, 32 refined poses and a crystal anchor;
the loader samples a bounded subset. Validation excludes crystal anchors.
The continuation uses Muon at $2\times10^{-4}$ plus AdamW at $3\times10^{-6}$,
effective global batch two, seed 45, 2,000-update warmup and cosine decay over
the final 50,000 of 100,000 updates, ending at 0.05 of the peak rates.
U70k was selected only on 1,035 internal complexes, scoring 622/1,035 (60.10\%)
Top-1 RMSD success, compared with 617/1,035 at U100k.

\subsection{Energy-based rigid refinement}
Refinement is a separate post-sampling operation with a fixed receptor:
\begin{equation}
 E=E_{\rm physical}+E_{\rm interaction},\qquad F_a=-\nabla_{x_a}E.
 \label{eq:effdock-energy}
\end{equation}
The physical terms are harmonic bond/angle restraints, periodic proper
torsions, planar/chiral impropers, softened Lennard--Jones interactions,
compact protein--ligand overlap barriers and repulsive receptor obstacles.
Physical parameters use EFF-FF version 2.2.0 and interaction parameters version
1.6.0. The chemical-constraint channel has scale zero, while the physical
chiral-improper scale is one. The input conformer's geometry, not the target
crystal pose, defines restraint reference values.
Bond restraints act on cut bonds, angles and impropers span fragments, and
proper-torsion weights sum to one per cut bond over its retained quadruplets.
\begin{align}
 E_{\rm bond}&=\tfrac12\sum_bk_b(d_b-d_b^0)^2,&
 E_{\rm angle}&=\tfrac12\sum_ak_a(\theta_a-\theta_a^0)^2,\\
 E_{\rm proper}&=\sum_pk_pw_p[1+\cos(n_p\phi_p-\delta_p)],\\
 E_{\rm improper}&=\tfrac12\sum_{i\in\rm planar}k_i[1-\cos(2\phi_i)]
 +\tfrac12\sum_{i\in\rm chiral}k_i\operatorname{wrap}(\phi_i-\phi_i^0)^2.
\end{align}
For $Q(u)=u^3(10-15u+6u^2)$, $S(d)=Q(\operatorname{clip}((8-d)/2,0,1))$
and $\rho=\sqrt{d^2+0.75^2}$,
\begin{equation}
 E_{{\rm LJ},ij}=s_{ij}S(d_{ij})D_{ij}
 [(X_{ij}/\rho_{ij})^{12}-2(X_{ij}/\rho_{ij})^6].
\end{equation}
Pair atom parameters use geometric mixing; the versioned table specifies
exclusions and the default 1--4 scale of 0.5. The steric term is
\begin{align}
 d_{ij}^{\rm safe}&=0.8(r_i^{\rm vdW}+r_j^{\rm vdW}),&
 p_{ij}&=0.1\log[1+e^{(d_{ij}^{\rm safe}-d_{ij})/0.1}],\\
 E_{\rm steric}&=10\sum_{ij}p_{ij}^2
 Q\!\left(\operatorname{clip}\left(\frac{d_{ij}^{\rm safe}+0.5-d_{ij}}{0.5},0,1\right)\right).
\end{align}
The interaction terms are hydrophobic, hydrogen bond, screened formal charge,
pi stacking, cation--pi, halogen bond and profile-specific metal coordination.
For bounded pair weights $w_{ij}$, site saturation is
\begin{equation}
 O_i=1-\prod_j[1-(1-10^{-7})w_{ij}],\qquad E_{\rm site}=-\epsilon\sum_iO_i.
\end{equation}
Symmetric ring-system saturation averages occupancies in both directions.
Screened formal-charge energy uses
\begin{equation}
 E_{\rm charge}=\sum_{ij}\frac{332.06371}{78.5}q_iq_j
 \frac{e^{-0.127\rho_{ij}}}{\rho_{ij}}S_q(d_{ij}^2),\quad
 \rho_{ij}=\sqrt{d_{ij}^2+0.75^2},
\end{equation}
with the quintic switch between squared distances $8^2$ and $10^2$.
The parameter tables and implementation specify all typed geometric gates,
metal profiles and repulsive-only obstacle terms. These diagnostic energies
are not calibrated binding free energies.

Energy-force projection uses atomic masses, unlike the learned head. For
fragment mass $M_f$, centre of mass $C_f$, centroid $T_f$ and $r_a=x_a-C_f$,
\begin{align}
 C_f&=\frac{\sum_{a\in f}m_ax_a}{M_f},&
 I_f&=\sum_{a\in f}m_a(\|r_a\|^2I_3-r_ar_a^{\mathsf T}),\\
 \Omega_f&=I_f^+\sum_{a\in f}r_a\times F_a,&
 V_f&=\frac{\sum_{a\in f}F_a}{M_f}+\Omega_f\times(T_f-C_f).
 \label{eq:effdock-energy-projection}
\end{align}
Independent per-pose backtracking tries $\alpha=0.5^b$, $b=0,\ldots,12$,
accepting finite energy no greater than
$E+10^{-10}\max(1,|E|)$. Each step is capped at $0.10\,\text{\AA}$
translation, five degrees rotation and $0.10\,\text{\AA}$ actual atom
movement. At most 100 iterations are run. Stop after 20 consecutive accepted
steps with maximum atom movement below $10^{-5}\,\text{\AA}$, or, from
iteration 25, five consecutive steps with energy decrease at most
$0.02+0.001\max(1,|E|)$. External runs use an $18\,\text{\AA}$ receptor
shell, $8\,\text{\AA}$ physical cutoff and $10\,\text{\AA}$ maximum
interaction cutoff. Training-bank refinement instead used displacement
threshold $0.01\,\text{\AA}$ with patience five. Failed line searches retain
the last accepted pose and an explicit status.

\subsection{Guided-sampling diagnostics}
The separately captioned guidance/budget and pocket/prior figures use
direct ODE guidance with $\eta=2$; the main condition has $\eta=0$.
Let $U(v,\omega)$ be the induced atom velocity and
$\operatorname{RMS}(U)=\sqrt{A^{-1}\sum_a\|U_a\|^2}$. Before capping,
\begin{equation}
 (\Delta v,\Delta\omega)=\eta\bar r_j
 \frac{\operatorname{RMS}(U(v_\theta,\omega_\theta))}
 {\operatorname{RMS}(U(V,\Omega))}(V,\Omega),\qquad
 \bar r_j=\frac{1}{\Delta t_j}\int_{t_j}^{t_{j+1}}\max(0,2t-1)\,dt.
\end{equation}
Here $(V,\Omega)$ projects energy forces capped at atom norm 20. The physical
LJ softcore decreases from 1.5 to $0.75\,\text{\AA}$ over the active
half-interval, evaluated at the active interval midpoint. The drift vanishes
when either RMS is at most $10^{-8}$ or the direction is nonfinite.
One scale per pose enforces added translation/angular velocity limits of
5 and 5 per unit flow time, and an estimated atom-displacement bound of
$0.25\,\text{\AA}$. The sampler adds the drift to the learned field before
its Euler update. Chemical-constraint drift is zero. This is separate from
the subsequent energy-based refinement.

\subsection{Selection, metrics and relatedness}
Predicted RMSD is $\widehat r_k=\max(0,\operatorname{expm1}
(\operatorname{clip}(\widehat y_k,-2,5)))$. Ordinary selection minimizes
$(\widehat r_k,k)$ lexicographically. Chirality filtering restricts this
minimum to candidates consistent with input tetrahedral stereochemistry,
falling back to ordinary selection if none passes. The main mask contains no
reference RMSD, PB test or E/Z criterion.

For $M$ complexes, selected index $k_i^*$, RMSD $r_{ik}$ and official
selected-pose PB validity $b_{ik}$,
\begin{equation}
 \mathrm{SR}=\frac{100}{M}\sum_i\mathbf1[r_{i,k_i^*}<2\,\text{\AA}],\quad
 \mathrm{SR}_{\rm PB}=\frac{100}{M}\sum_i\mathbf1[r_{i,k_i^*}<2\,\text{\AA}]b_{i,k_i^*}.
\end{equation}
RMSD uses symmetry-aware no-alignment heavy-atom CalcRMS. Solid bars show
$\mathrm{SR}_{\rm PB}$; hatched extensions show $\mathrm{SR}-\mathrm{SR}_{\rm PB}$.
Top-$k$ uses the full-bank confidence order, while generation-prefix coverage
uses saved generation indices. Both reach the RMSD-only oracle at 100; prefix
reselection is a distinct, potentially nonmonotone selected-pose endpoint.

The external cohorts are Astex Diverse Set (85), PoseBusters v2 (308),
PhiBench (206), FoldBench (558) and auxiliary OpenBind (925). PhiBench includes
three reconstructed cases; OpenBind includes two flagged noncovalent
approximations of covalent systems. FoldBench PB evaluation includes three
energy-reference InChI compatibility-repair shards. Original denominators are
retained. Report three-seed means and sample SD, except paired bootstrap
figures: average paired differences over seeds within each complex, then
resample complexes or exact-PDB groups 2,000 times and report percentile
95\% CIs in percentage points. Group resampling retains member complexes
and uses their complex-weighted mean; it is not protein-family clustering.

Sequence relatedness uses the 47,277 docking samples. For query binding
chains $Q$ and a training sample $h$, define
\begin{equation}
 s(Q,h)=100\frac{\max_{\pi}\sum_{q\in Q}n_{\rm identical}(q,\pi(q))}
 {\sum_{q\in Q}|q|},\qquad s_{\max}(Q)=\max_{h\in\mathcal H}s(Q,h),
\end{equation}
where $\pi$ is a one-to-one assignment to eligible binding chains within
that training sample. Unaligned residues, unmatched chains and unknown
residues do not contribute identical-residue evidence; all observed query
residues remain in the denominator. Ligand--protein overlap requires exact
canonical heavy-isomeric ligand identity and $s(Q,h)\ge70$ in the same
training sample, searching all exact-ligand witnesses. The threshold is
descriptive, not a certified pocket match or binary leakage definition.
Six-cohort composition adds Validation (1,076), without implying a matched
three-seed validation performance bank. Opened external results are
descriptive and do not select checkpoints or tune settings.
.
% Implementation references and full typed-energy parameter tables are in docs/methods.
\section{Methods}
\label{sec:effdock-methods}

\subsection{Task, data and information boundary}
EFF-Dock predicts ligand poses given a receptor, ligand chemistry and a
supplied pocket centre. The experiments are retrospective supplied-pocket
redocking: pocket centres may be reference-ligand-derived and receptors may
be holo structures. The reference pose supplies supervision, mapping and
evaluation labels; its coordinates do not enter the learned vector field,
confidence ranker or minimized refinement objective.

Training structures originate from PLINDER 2024-06/v2. A sample is a system
and ligand-instance pair. The preserved split contains 47,310 training and
1,076 validation IDs. The released docking fine-tune loader contains 47,277
samples; confidence uses 43,092 training and 1,035 validation samples. Their
training intersection is 43,067, with 4,210 docking-only and 25 confidence-only
samples. These independently filtered inventories are not nested and do not
represent the union of every predecessor checkpoint's training exposure.
Public membership records provide IDs, exclusions and SHA-256 hashes.

The historical external exclusion combined a ligand-match condition with a
pocket-community-match condition whose reference communities were identified
from benchmark PDB hits in the processed pool. It is distinct from the newer
strict split builder. No post-hoc exclusion changes the released benchmarks.

\subsection{Fragment representation and equivariant model}
Ligands are partitioned at eligible single, non-ring, non-planar-conjugated
rotatable bonds, retaining at least two heavy atoms per fragment after
singleton merging. Let $f(a)$ assign atom $a$ to fragment $f$, $T_f$ be its
centroid, $q_f$ a unit quaternion and $\ell_a$ its fixed local coordinate:
\begin{equation}
 T_f=\frac{1}{|f|}\sum_{a\in f}x_a,\qquad
 x_a(T,q)=R(q_{f(a)})\ell_a+T_{f(a)}.
 \label{eq:effdock-reconstruction}
\end{equation}
The heterogeneous graph contains ligand atoms, fragments, protein atoms and
residue virtual nodes. Static topology and dynamic protein--ligand contacts
within $5\,\text{\AA}$ condition six equivariant interaction layers, using
32 radial basis functions and spherical harmonics through degree two. Hidden
irreps are $384\times0e+32\times1o+32\times1e+16\times2e+16\times2o$,
with flattened width 736. Irrep notation alone does not establish full
reflection equivariance of stereochemical inputs; the pose contract uses
proper rotations and joint SE(3) transformations.

The model outputs atom vector fields $g_a$. Define $r_a=x_a-T_{f(a)}$.
Unit-weight Newton--Euler aggregation gives
\begin{align}
 v_f&=\frac{1}{|f|}\sum_{a\in f}g_a,&
 \tau_f&=\sum_{a\in f}r_a\times g_a,\\
 I_f&=\sum_{a\in f}(\|r_a\|^2I_3-r_ar_a^{\mathsf T}),&
 \omega_f&=I_f^+\tau_f.
 \label{eq:effdock-head}
\end{align}
The pseudoinverse retains inertia eigenvalues strictly above 1\% of the
largest eigenvalue, yielding an observable-rotation projector $P_f$.

\subsection{Conditional flow and training objective}
For supplied centre $c$, translation priors are
$T_f^0=c+\sigma\epsilon_f$, $\epsilon_f\sim\mathcal N(0,I_3)$;
initial orientations are uniform on SO(3). Conditional interpolation is
\begin{align}
 T_f^t&=(1-t)T_f^0+tT_f^1,&
 q_f^t&=\operatorname{SLERP}(q_f^0,q_f^1;t),\\
 v_f^*&=T_f^1-T_f^0,&
 \omega_f^*&=\operatorname{axisAngle}(q_f^1\otimes(q_f^0)^{-1}).
 \label{eq:effdock-flow-targets}
\end{align}
Quaternion updates use left multiplication. With $F$ fragments and
$F_\omega$ fragments having more than one atom, the flow terms are
\begin{equation}
 L_v=\frac{1}{3F}\sum_f\|v_f-v_f^*\|^2,\qquad
 L_\omega=\frac{1}{3F_\omega}\sum_{f:|f|>1}
 \|\omega_f-P_f\omega_f^*\|^2.
 \label{eq:effdock-flow-loss}
\end{equation}
A zero eligible count yields zero angular loss. The angular denominator is
not the number of observable eigenvectors. The released fragmenter avoids
singletons. Optional time weighting of these flow terms is disabled.

The atom auxiliary loss compares instantaneous velocities, not endpoint
coordinates. For $r_a^t=x_a^t-T_{f(a)}^t$,
\begin{equation}
 u_a=v_{f(a)}+\omega_{f(a)}\times r_a^t,\quad
 u_a^*=v_{f(a)}^*+\omega_{f(a)}^*\times r_a^t,\quad
 L_{\rm atom}=\frac{1}{3A}\sum_a\|u_a-u_a^*\|^2.
 \label{eq:effdock-atom-loss}
\end{equation}
The auxiliary target uses the unprojected angular field; an exactly
unobservable rotation induces zero atom motion.

A separate distance-geometry term uses the one-step endpoint
\begin{equation}
 \widehat T_f^1=T_f^t+(1-t)v_f,\qquad
 \widehat q_f^1=\operatorname{Exp}_q((1-t)\omega_f)\otimes q_f^t.
\end{equation}
Here $\operatorname{Exp}_q$ converts an axis-angle vector to a quaternion.
For unordered inter-fragment atom pairs $\mathcal P_b$ in complex $b$,
\begin{align}
 D_b&=\frac{1}{|\mathcal P_b|}\sum_{(a,c)\in\mathcal P_b}
 \left(\|\widehat x_a^1-\widehat x_c^1\|-\|x_a^1-x_c^1\|\right)^2,\\
 L_{\rm DG}&=\frac{\sum_{b:|\mathcal P_b|>0}t_b^2D_b}
 {\sum_{b:|\mathcal P_b|>0}t_b^2},\qquad
 L=L_v+8L_\omega+0.3L_{\rm atom}+3L_{\rm DG}.
 \label{eq:effdock-total-loss}
\end{align}
A zero denominator yields zero $L_{\rm DG}$. Distributed reductions reproduce
global fragment, atom and weighted-complex means for the respective terms;
these are not equal-weight-per-complex reductions throughout.

The released 50k fine-tune samples
\begin{align}
 p(t)&=0.80p_{\rm SF}(t)+0.10U(0,0.3)+0.10\delta_0,\\
 p_{\rm SF}&=0.02U(0,1)+0.98\operatorname{Law}
 [\operatorname{sigmoid}(0.8+1.7Z)],\quad Z\sim\mathcal N(0,1).
\end{align}
It uses fresh AdamW state, global batch 64 on four GPUs, peak learning rate
$2\times10^{-5}$, weight decay 0.01, gradient clipping 1, and EMA decay 0.999.
Warmup lasts 1,000 updates; cosine decay begins at 40,000 and reaches
$2\times10^{-6}$ at 50,000. The released docking state is terminal EMA.

\subsection{Sampling and confidence}
Main sampling uses $N=100$ candidates, $S=10$ steps, $\sigma=2\,\text{\AA}$
and a $10\,\text{\AA}$ crop. The late-power-three Euler grid and updates are
\begin{align}
 t_j&=1-(1-j/S)^3,\quad j=0,\ldots,S,\\
 T_f^{j+1}&=T_f^j+\Delta t_jv_{\theta,f}(t_j),\qquad
 q_f^{j+1}=\operatorname{Exp}_q(\Delta t_j\omega_{\theta,f}(t_j))\otimes q_f^j.
 \label{eq:effdock-sampler}
\end{align}
Raw and refined banks preserve generation indices and are independently scored.

The confidence scorer uses four equivariant layers, paired docking features,
and attention pooling over global and contact features. It predicts log-one-plus pose
RMSD $\widehat y_k$, pose-success logit $z_k$, atom log displacement
$\widehat y_{ka}$ and atom-success logit $z_{ka}$. Atom and pose regression
use Huber loss with transition one; classification uses BCE with logits and
strict $2\,\text{\AA}$ labels. Let $y_k=\log(1+r_k)$:
\begin{align}
 L_{\rm rank}&=\underset{y_j-y_i>0.3}{\operatorname{mean}}
 \max(0,0.05-(\widehat y_j-\widehat y_i)),\\
 q_k&=\operatorname{sigmoid}((2-r_k)/0.2),\quad
 p_k=\frac{q_k}{\sum_jq_j},\qquad
 L_{\rm list}=-\sum_kp_k\log\operatorname{softmax}(z)_k,\\
 L_{\rm conf}&=0.2L_{\rm atom}+0.2L_{\rm atom\mbox{-}ok}
 +0.3L_{\rm pose}+0.4L_{\rm pose\mbox{-}ok}
 +0.1L_{\rm rank}+L_{\rm list}.
 \label{eq:effdock-confidence}
\end{align}
Here the confidence atom term is displacement regression, distinct from the
docking atom-velocity term in Eq.~\ref{eq:effdock-atom-loss}. Empty eligible
pairs give zero rank loss. Numerically zero quality sums use a uniform target;
one-candidate pools have zero listwise loss.

Training banks contain 32 raw poses, 32 refined poses and a crystal anchor;
the loader samples a bounded subset. Validation excludes crystal anchors.
The continuation uses Muon at $2\times10^{-4}$ plus AdamW at $3\times10^{-6}$,
effective global batch two, seed 45, 2,000-update warmup and cosine decay over
the final 50,000 of 100,000 updates, ending at 0.05 of the peak rates.
U70k was selected only on 1,035 internal complexes, scoring 622/1,035 (60.10\%)
Top-1 RMSD success, compared with 617/1,035 at U100k.

\subsection{Energy-based rigid refinement}
Refinement is a separate post-sampling operation with a fixed receptor:
\begin{equation}
 E=E_{\rm physical}+E_{\rm interaction},\qquad F_a=-\nabla_{x_a}E.
 \label{eq:effdock-energy}
\end{equation}
The physical terms are harmonic bond/angle restraints, periodic proper
torsions, planar/chiral impropers, softened Lennard--Jones interactions,
compact protein--ligand overlap barriers and repulsive receptor obstacles.
Physical parameters use EFF-FF version 2.2.0 and interaction parameters version
1.6.0. The chemical-constraint channel has scale zero, while the physical
chiral-improper scale is one. The input conformer's geometry, not the target
crystal pose, defines restraint reference values.
Bond restraints act on cut bonds, angles and impropers span fragments, and
proper-torsion weights sum to one per cut bond over its retained quadruplets.
\begin{align}
 E_{\rm bond}&=\tfrac12\sum_bk_b(d_b-d_b^0)^2,&
 E_{\rm angle}&=\tfrac12\sum_ak_a(\theta_a-\theta_a^0)^2,\\
 E_{\rm proper}&=\sum_pk_pw_p[1+\cos(n_p\phi_p-\delta_p)],\\
 E_{\rm improper}&=\tfrac12\sum_{i\in\rm planar}k_i[1-\cos(2\phi_i)]
 +\tfrac12\sum_{i\in\rm chiral}k_i\operatorname{wrap}(\phi_i-\phi_i^0)^2.
\end{align}
For $Q(u)=u^3(10-15u+6u^2)$, $S(d)=Q(\operatorname{clip}((8-d)/2,0,1))$
and $\rho=\sqrt{d^2+0.75^2}$,
\begin{equation}
 E_{{\rm LJ},ij}=s_{ij}S(d_{ij})D_{ij}
 [(X_{ij}/\rho_{ij})^{12}-2(X_{ij}/\rho_{ij})^6].
\end{equation}
Pair atom parameters use geometric mixing; the versioned table specifies
exclusions and the default 1--4 scale of 0.5. The steric term is
\begin{align}
 d_{ij}^{\rm safe}&=0.8(r_i^{\rm vdW}+r_j^{\rm vdW}),&
 p_{ij}&=0.1\log[1+e^{(d_{ij}^{\rm safe}-d_{ij})/0.1}],\\
 E_{\rm steric}&=10\sum_{ij}p_{ij}^2
 Q\!\left(\operatorname{clip}\left(\frac{d_{ij}^{\rm safe}+0.5-d_{ij}}{0.5},0,1\right)\right).
\end{align}
The interaction terms are hydrophobic, hydrogen bond, screened formal charge,
pi stacking, cation--pi, halogen bond and profile-specific metal coordination.
For bounded pair weights $w_{ij}$, site saturation is
\begin{equation}
 O_i=1-\prod_j[1-(1-10^{-7})w_{ij}],\qquad E_{\rm site}=-\epsilon\sum_iO_i.
\end{equation}
Symmetric ring-system saturation averages occupancies in both directions.
Screened formal-charge energy uses
\begin{equation}
 E_{\rm charge}=\sum_{ij}\frac{332.06371}{78.5}q_iq_j
 \frac{e^{-0.127\rho_{ij}}}{\rho_{ij}}S_q(d_{ij}^2),\quad
 \rho_{ij}=\sqrt{d_{ij}^2+0.75^2},
\end{equation}
with the quintic switch between squared distances $8^2$ and $10^2$.
The parameter tables and implementation specify all typed geometric gates,
metal profiles and repulsive-only obstacle terms. These diagnostic energies
are not calibrated binding free energies.

Energy-force projection uses atomic masses, unlike the learned head. For
fragment mass $M_f$, centre of mass $C_f$, centroid $T_f$ and $r_a=x_a-C_f$,
\begin{align}
 C_f&=\frac{\sum_{a\in f}m_ax_a}{M_f},&
 I_f&=\sum_{a\in f}m_a(\|r_a\|^2I_3-r_ar_a^{\mathsf T}),\\
 \Omega_f&=I_f^+\sum_{a\in f}r_a\times F_a,&
 V_f&=\frac{\sum_{a\in f}F_a}{M_f}+\Omega_f\times(T_f-C_f).
 \label{eq:effdock-energy-projection}
\end{align}
Independent per-pose backtracking tries $\alpha=0.5^b$, $b=0,\ldots,12$,
accepting finite energy no greater than
$E+10^{-10}\max(1,|E|)$. Each step is capped at $0.10\,\text{\AA}$
translation, five degrees rotation and $0.10\,\text{\AA}$ actual atom
movement. At most 100 iterations are run. Stop after 20 consecutive accepted
steps with maximum atom movement below $10^{-5}\,\text{\AA}$, or, from
iteration 25, five consecutive steps with energy decrease at most
$0.02+0.001\max(1,|E|)$. External runs use an $18\,\text{\AA}$ receptor
shell, $8\,\text{\AA}$ physical cutoff and $10\,\text{\AA}$ maximum
interaction cutoff. Training-bank refinement instead used displacement
threshold $0.01\,\text{\AA}$ with patience five. Failed line searches retain
the last accepted pose and an explicit status.

\subsection{Guided-sampling diagnostics}
The separately captioned guidance/budget and pocket/prior figures use
direct ODE guidance with $\eta=2$; the main condition has $\eta=0$.
Let $U(v,\omega)$ be the induced atom velocity and
$\operatorname{RMS}(U)=\sqrt{A^{-1}\sum_a\|U_a\|^2}$. Before capping,
\begin{equation}
 (\Delta v,\Delta\omega)=\eta\bar r_j
 \frac{\operatorname{RMS}(U(v_\theta,\omega_\theta))}
 {\operatorname{RMS}(U(V,\Omega))}(V,\Omega),\qquad
 \bar r_j=\frac{1}{\Delta t_j}\int_{t_j}^{t_{j+1}}\max(0,2t-1)\,dt.
\end{equation}
Here $(V,\Omega)$ projects energy forces capped at atom norm 20. The physical
LJ softcore decreases from 1.5 to $0.75\,\text{\AA}$ over the active
half-interval, evaluated at the active interval midpoint. The drift vanishes
when either RMS is at most $10^{-8}$ or the direction is nonfinite.
One scale per pose enforces added translation/angular velocity limits of
5 and 5 per unit flow time, and an estimated atom-displacement bound of
$0.25\,\text{\AA}$. The sampler adds the drift to the learned field before
its Euler update. Chemical-constraint drift is zero. This is separate from
the subsequent energy-based refinement.

\subsection{Selection, metrics and relatedness}
Predicted RMSD is $\widehat r_k=\max(0,\operatorname{expm1}
(\operatorname{clip}(\widehat y_k,-2,5)))$. Ordinary selection minimizes
$(\widehat r_k,k)$ lexicographically. Chirality filtering restricts this
minimum to candidates consistent with input tetrahedral stereochemistry,
falling back to ordinary selection if none passes. The main mask contains no
reference RMSD, PB test or E/Z criterion.

For $M$ complexes, selected index $k_i^*$, RMSD $r_{ik}$ and official
selected-pose PB validity $b_{ik}$,
\begin{equation}
 \mathrm{SR}=\frac{100}{M}\sum_i\mathbf1[r_{i,k_i^*}<2\,\text{\AA}],\quad
 \mathrm{SR}_{\rm PB}=\frac{100}{M}\sum_i\mathbf1[r_{i,k_i^*}<2\,\text{\AA}]b_{i,k_i^*}.
\end{equation}
RMSD uses symmetry-aware no-alignment heavy-atom CalcRMS. Solid bars show
$\mathrm{SR}_{\rm PB}$; hatched extensions show $\mathrm{SR}-\mathrm{SR}_{\rm PB}$.
Top-$k$ uses the full-bank confidence order, while generation-prefix coverage
uses saved generation indices. Both reach the RMSD-only oracle at 100; prefix
reselection is a distinct, potentially nonmonotone selected-pose endpoint.

The external cohorts are Astex Diverse Set (85), PoseBusters v2 (308),
PhiBench (206), FoldBench (558) and auxiliary OpenBind (925). PhiBench includes
three reconstructed cases; OpenBind includes two flagged noncovalent
approximations of covalent systems. FoldBench PB evaluation includes three
energy-reference InChI compatibility-repair shards. Original denominators are
retained. Report three-seed means and sample SD, except paired bootstrap
figures: average paired differences over seeds within each complex, then
resample complexes or exact-PDB groups 2,000 times and report percentile
95\% CIs in percentage points. Group resampling retains member complexes
and uses their complex-weighted mean; it is not protein-family clustering.

Sequence relatedness uses the 47,277 docking samples. For query binding
chains $Q$ and a training sample $h$, define
\begin{equation}
 s(Q,h)=100\frac{\max_{\pi}\sum_{q\in Q}n_{\rm identical}(q,\pi(q))}
 {\sum_{q\in Q}|q|},\qquad s_{\max}(Q)=\max_{h\in\mathcal H}s(Q,h),
\end{equation}
where $\pi$ is a one-to-one assignment to eligible binding chains within
that training sample. Unaligned residues, unmatched chains and unknown
residues do not contribute identical-residue evidence; all observed query
residues remain in the denominator. Ligand--protein overlap requires exact
canonical heavy-isomeric ligand identity and $s(Q,h)\ge70$ in the same
training sample, searching all exact-ligand witnesses. The threshold is
descriptive, not a certified pocket match or binary leakage definition.
Six-cohort composition adds Validation (1,076), without implying a matched
three-seed validation performance bank. Opened external results are
descriptive and do not select checkpoints or tune settings.
.
% Implementation references and full typed-energy parameter tables are in docs/methods.
\section{Methods}
\label{sec:effdock-methods}

\subsection{Task, data and information boundary}
EFF-Dock predicts ligand poses given a receptor, ligand chemistry and a
supplied pocket centre. The experiments are retrospective supplied-pocket
redocking: pocket centres may be reference-ligand-derived and receptors may
be holo structures. The reference pose supplies supervision, mapping and
evaluation labels; its coordinates do not enter the learned vector field,
confidence ranker or minimized refinement objective.

Training structures originate from PLINDER 2024-06/v2. A sample is a system
and ligand-instance pair. The preserved split contains 47,310 training and
1,076 validation IDs. The released docking fine-tune loader contains 47,277
samples; confidence uses 43,092 training and 1,035 validation samples. Their
training intersection is 43,067, with 4,210 docking-only and 25 confidence-only
samples. These independently filtered inventories are not nested and do not
represent the union of every predecessor checkpoint's training exposure.
Public membership records provide IDs, exclusions and SHA-256 hashes.

The historical external exclusion combined a ligand-match condition with a
pocket-community-match condition whose reference communities were identified
from benchmark PDB hits in the processed pool. It is distinct from the newer
strict split builder. No post-hoc exclusion changes the released benchmarks.

\subsection{Fragment representation and equivariant model}
Ligands are partitioned at eligible single, non-ring, non-planar-conjugated
rotatable bonds, retaining at least two heavy atoms per fragment after
singleton merging. Let $f(a)$ assign atom $a$ to fragment $f$, $T_f$ be its
centroid, $q_f$ a unit quaternion and $\ell_a$ its fixed local coordinate:
\begin{equation}
 T_f=\frac{1}{|f|}\sum_{a\in f}x_a,\qquad
 x_a(T,q)=R(q_{f(a)})\ell_a+T_{f(a)}.
 \label{eq:effdock-reconstruction}
\end{equation}
The heterogeneous graph contains ligand atoms, fragments, protein atoms and
residue virtual nodes. Static topology and dynamic protein--ligand contacts
within $5\,\text{\AA}$ condition six equivariant interaction layers, using
32 radial basis functions and spherical harmonics through degree two. Hidden
irreps are $384\times0e+32\times1o+32\times1e+16\times2e+16\times2o$,
with flattened width 736. Irrep notation alone does not establish full
reflection equivariance of stereochemical inputs; the pose contract uses
proper rotations and joint SE(3) transformations.

The model outputs atom vector fields $g_a$. Define $r_a=x_a-T_{f(a)}$.
Unit-weight Newton--Euler aggregation gives
\begin{align}
 v_f&=\frac{1}{|f|}\sum_{a\in f}g_a,&
 \tau_f&=\sum_{a\in f}r_a\times g_a,\\
 I_f&=\sum_{a\in f}(\|r_a\|^2I_3-r_ar_a^{\mathsf T}),&
 \omega_f&=I_f^+\tau_f.
 \label{eq:effdock-head}
\end{align}
The pseudoinverse retains inertia eigenvalues strictly above 1\% of the
largest eigenvalue, yielding an observable-rotation projector $P_f$.

\subsection{Conditional flow and training objective}
For supplied centre $c$, translation priors are
$T_f^0=c+\sigma\epsilon_f$, $\epsilon_f\sim\mathcal N(0,I_3)$;
initial orientations are uniform on SO(3). Conditional interpolation is
\begin{align}
 T_f^t&=(1-t)T_f^0+tT_f^1,&
 q_f^t&=\operatorname{SLERP}(q_f^0,q_f^1;t),\\
 v_f^*&=T_f^1-T_f^0,&
 \omega_f^*&=\operatorname{axisAngle}(q_f^1\otimes(q_f^0)^{-1}).
 \label{eq:effdock-flow-targets}
\end{align}
Quaternion updates use left multiplication. With $F$ fragments and
$F_\omega$ fragments having more than one atom, the flow terms are
\begin{equation}
 L_v=\frac{1}{3F}\sum_f\|v_f-v_f^*\|^2,\qquad
 L_\omega=\frac{1}{3F_\omega}\sum_{f:|f|>1}
 \|\omega_f-P_f\omega_f^*\|^2.
 \label{eq:effdock-flow-loss}
\end{equation}
A zero eligible count yields zero angular loss. The angular denominator is
not the number of observable eigenvectors. The released fragmenter avoids
singletons. Optional time weighting of these flow terms is disabled.

The atom auxiliary loss compares instantaneous velocities, not endpoint
coordinates. For $r_a^t=x_a^t-T_{f(a)}^t$,
\begin{equation}
 u_a=v_{f(a)}+\omega_{f(a)}\times r_a^t,\quad
 u_a^*=v_{f(a)}^*+\omega_{f(a)}^*\times r_a^t,\quad
 L_{\rm atom}=\frac{1}{3A}\sum_a\|u_a-u_a^*\|^2.
 \label{eq:effdock-atom-loss}
\end{equation}
The auxiliary target uses the unprojected angular field; an exactly
unobservable rotation induces zero atom motion.

A separate distance-geometry term uses the one-step endpoint
\begin{equation}
 \widehat T_f^1=T_f^t+(1-t)v_f,\qquad
 \widehat q_f^1=\operatorname{Exp}_q((1-t)\omega_f)\otimes q_f^t.
\end{equation}
Here $\operatorname{Exp}_q$ converts an axis-angle vector to a quaternion.
For unordered inter-fragment atom pairs $\mathcal P_b$ in complex $b$,
\begin{align}
 D_b&=\frac{1}{|\mathcal P_b|}\sum_{(a,c)\in\mathcal P_b}
 \left(\|\widehat x_a^1-\widehat x_c^1\|-\|x_a^1-x_c^1\|\right)^2,\\
 L_{\rm DG}&=\frac{\sum_{b:|\mathcal P_b|>0}t_b^2D_b}
 {\sum_{b:|\mathcal P_b|>0}t_b^2},\qquad
 L=L_v+8L_\omega+0.3L_{\rm atom}+3L_{\rm DG}.
 \label{eq:effdock-total-loss}
\end{align}
A zero denominator yields zero $L_{\rm DG}$. Distributed reductions reproduce
global fragment, atom and weighted-complex means for the respective terms;
these are not equal-weight-per-complex reductions throughout.

The released 50k fine-tune samples
\begin{align}
 p(t)&=0.80p_{\rm SF}(t)+0.10U(0,0.3)+0.10\delta_0,\\
 p_{\rm SF}&=0.02U(0,1)+0.98\operatorname{Law}
 [\operatorname{sigmoid}(0.8+1.7Z)],\quad Z\sim\mathcal N(0,1).
\end{align}
It uses fresh AdamW state, global batch 64 on four GPUs, peak learning rate
$2\times10^{-5}$, weight decay 0.01, gradient clipping 1, and EMA decay 0.999.
Warmup lasts 1,000 updates; cosine decay begins at 40,000 and reaches
$2\times10^{-6}$ at 50,000. The released docking state is terminal EMA.

\subsection{Sampling and confidence}
Main sampling uses $N=100$ candidates, $S=10$ steps, $\sigma=2\,\text{\AA}$
and a $10\,\text{\AA}$ crop. The late-power-three Euler grid and updates are
\begin{align}
 t_j&=1-(1-j/S)^3,\quad j=0,\ldots,S,\\
 T_f^{j+1}&=T_f^j+\Delta t_jv_{\theta,f}(t_j),\qquad
 q_f^{j+1}=\operatorname{Exp}_q(\Delta t_j\omega_{\theta,f}(t_j))\otimes q_f^j.
 \label{eq:effdock-sampler}
\end{align}
Raw and refined banks preserve generation indices and are independently scored.

The confidence scorer uses four equivariant layers, paired docking features,
and attention pooling over global and contact features. It predicts log-one-plus pose
RMSD $\widehat y_k$, pose-success logit $z_k$, atom log displacement
$\widehat y_{ka}$ and atom-success logit $z_{ka}$. Atom and pose regression
use Huber loss with transition one; classification uses BCE with logits and
strict $2\,\text{\AA}$ labels. Let $y_k=\log(1+r_k)$:
\begin{align}
 L_{\rm rank}&=\underset{y_j-y_i>0.3}{\operatorname{mean}}
 \max(0,0.05-(\widehat y_j-\widehat y_i)),\\
 q_k&=\operatorname{sigmoid}((2-r_k)/0.2),\quad
 p_k=\frac{q_k}{\sum_jq_j},\qquad
 L_{\rm list}=-\sum_kp_k\log\operatorname{softmax}(z)_k,\\
 L_{\rm conf}&=0.2L_{\rm atom}+0.2L_{\rm atom\mbox{-}ok}
 +0.3L_{\rm pose}+0.4L_{\rm pose\mbox{-}ok}
 +0.1L_{\rm rank}+L_{\rm list}.
 \label{eq:effdock-confidence}
\end{align}
Here the confidence atom term is displacement regression, distinct from the
docking atom-velocity term in Eq.~\ref{eq:effdock-atom-loss}. Empty eligible
pairs give zero rank loss. Numerically zero quality sums use a uniform target;
one-candidate pools have zero listwise loss.

Training banks contain 32 raw poses, 32 refined poses and a crystal anchor;
the loader samples a bounded subset. Validation excludes crystal anchors.
The continuation uses Muon at $2\times10^{-4}$ plus AdamW at $3\times10^{-6}$,
effective global batch two, seed 45, 2,000-update warmup and cosine decay over
the final 50,000 of 100,000 updates, ending at 0.05 of the peak rates.
U70k was selected only on 1,035 internal complexes, scoring 622/1,035 (60.10\%)
Top-1 RMSD success, compared with 617/1,035 at U100k.

\subsection{Energy-based rigid refinement}
Refinement is a separate post-sampling operation with a fixed receptor:
\begin{equation}
 E=E_{\rm physical}+E_{\rm interaction},\qquad F_a=-\nabla_{x_a}E.
 \label{eq:effdock-energy}
\end{equation}
The physical terms are harmonic bond/angle restraints, periodic proper
torsions, planar/chiral impropers, softened Lennard--Jones interactions,
compact protein--ligand overlap barriers and repulsive receptor obstacles.
Physical parameters use EFF-FF version 2.2.0 and interaction parameters version
1.6.0. The chemical-constraint channel has scale zero, while the physical
chiral-improper scale is one. The input conformer's geometry, not the target
crystal pose, defines restraint reference values.
Bond restraints act on cut bonds, angles and impropers span fragments, and
proper-torsion weights sum to one per cut bond over its retained quadruplets.
\begin{align}
 E_{\rm bond}&=\tfrac12\sum_bk_b(d_b-d_b^0)^2,&
 E_{\rm angle}&=\tfrac12\sum_ak_a(\theta_a-\theta_a^0)^2,\\
 E_{\rm proper}&=\sum_pk_pw_p[1+\cos(n_p\phi_p-\delta_p)],\\
 E_{\rm improper}&=\tfrac12\sum_{i\in\rm planar}k_i[1-\cos(2\phi_i)]
 +\tfrac12\sum_{i\in\rm chiral}k_i\operatorname{wrap}(\phi_i-\phi_i^0)^2.
\end{align}
For $Q(u)=u^3(10-15u+6u^2)$, $S(d)=Q(\operatorname{clip}((8-d)/2,0,1))$
and $\rho=\sqrt{d^2+0.75^2}$,
\begin{equation}
 E_{{\rm LJ},ij}=s_{ij}S(d_{ij})D_{ij}
 [(X_{ij}/\rho_{ij})^{12}-2(X_{ij}/\rho_{ij})^6].
\end{equation}
Pair atom parameters use geometric mixing; the versioned table specifies
exclusions and the default 1--4 scale of 0.5. The steric term is
\begin{align}
 d_{ij}^{\rm safe}&=0.8(r_i^{\rm vdW}+r_j^{\rm vdW}),&
 p_{ij}&=0.1\log[1+e^{(d_{ij}^{\rm safe}-d_{ij})/0.1}],\\
 E_{\rm steric}&=10\sum_{ij}p_{ij}^2
 Q\!\left(\operatorname{clip}\left(\frac{d_{ij}^{\rm safe}+0.5-d_{ij}}{0.5},0,1\right)\right).
\end{align}
The interaction terms are hydrophobic, hydrogen bond, screened formal charge,
pi stacking, cation--pi, halogen bond and profile-specific metal coordination.
For bounded pair weights $w_{ij}$, site saturation is
\begin{equation}
 O_i=1-\prod_j[1-(1-10^{-7})w_{ij}],\qquad E_{\rm site}=-\epsilon\sum_iO_i.
\end{equation}
Symmetric ring-system saturation averages occupancies in both directions.
Screened formal-charge energy uses
\begin{equation}
 E_{\rm charge}=\sum_{ij}\frac{332.06371}{78.5}q_iq_j
 \frac{e^{-0.127\rho_{ij}}}{\rho_{ij}}S_q(d_{ij}^2),\quad
 \rho_{ij}=\sqrt{d_{ij}^2+0.75^2},
\end{equation}
with the quintic switch between squared distances $8^2$ and $10^2$.
The parameter tables and implementation specify all typed geometric gates,
metal profiles and repulsive-only obstacle terms. These diagnostic energies
are not calibrated binding free energies.

Energy-force projection uses atomic masses, unlike the learned head. For
fragment mass $M_f$, centre of mass $C_f$, centroid $T_f$ and $r_a=x_a-C_f$,
\begin{align}
 C_f&=\frac{\sum_{a\in f}m_ax_a}{M_f},&
 I_f&=\sum_{a\in f}m_a(\|r_a\|^2I_3-r_ar_a^{\mathsf T}),\\
 \Omega_f&=I_f^+\sum_{a\in f}r_a\times F_a,&
 V_f&=\frac{\sum_{a\in f}F_a}{M_f}+\Omega_f\times(T_f-C_f).
 \label{eq:effdock-energy-projection}
\end{align}
Independent per-pose backtracking tries $\alpha=0.5^b$, $b=0,\ldots,12$,
accepting finite energy no greater than
$E+10^{-10}\max(1,|E|)$. Each step is capped at $0.10\,\text{\AA}$
translation, five degrees rotation and $0.10\,\text{\AA}$ actual atom
movement. At most 100 iterations are run. Stop after 20 consecutive accepted
steps with maximum atom movement below $10^{-5}\,\text{\AA}$, or, from
iteration 25, five consecutive steps with energy decrease at most
$0.02+0.001\max(1,|E|)$. External runs use an $18\,\text{\AA}$ receptor
shell, $8\,\text{\AA}$ physical cutoff and $10\,\text{\AA}$ maximum
interaction cutoff. Training-bank refinement instead used displacement
threshold $0.01\,\text{\AA}$ with patience five. Failed line searches retain
the last accepted pose and an explicit status.

\subsection{Guided-sampling diagnostics}
The separately captioned guidance/budget and pocket/prior figures use
direct ODE guidance with $\eta=2$; the main condition has $\eta=0$.
Let $U(v,\omega)$ be the induced atom velocity and
$\operatorname{RMS}(U)=\sqrt{A^{-1}\sum_a\|U_a\|^2}$. Before capping,
\begin{equation}
 (\Delta v,\Delta\omega)=\eta\bar r_j
 \frac{\operatorname{RMS}(U(v_\theta,\omega_\theta))}
 {\operatorname{RMS}(U(V,\Omega))}(V,\Omega),\qquad
 \bar r_j=\frac{1}{\Delta t_j}\int_{t_j}^{t_{j+1}}\max(0,2t-1)\,dt.
\end{equation}
Here $(V,\Omega)$ projects energy forces capped at atom norm 20. The physical
LJ softcore decreases from 1.5 to $0.75\,\text{\AA}$ over the active
half-interval, evaluated at the active interval midpoint. The drift vanishes
when either RMS is at most $10^{-8}$ or the direction is nonfinite.
One scale per pose enforces added translation/angular velocity limits of
5 and 5 per unit flow time, and an estimated atom-displacement bound of
$0.25\,\text{\AA}$. The sampler adds the drift to the learned field before
its Euler update. Chemical-constraint drift is zero. This is separate from
the subsequent energy-based refinement.

\subsection{Selection, metrics and relatedness}
Predicted RMSD is $\widehat r_k=\max(0,\operatorname{expm1}
(\operatorname{clip}(\widehat y_k,-2,5)))$. Ordinary selection minimizes
$(\widehat r_k,k)$ lexicographically. Chirality filtering restricts this
minimum to candidates consistent with input tetrahedral stereochemistry,
falling back to ordinary selection if none passes. The main mask contains no
reference RMSD, PB test or E/Z criterion.

For $M$ complexes, selected index $k_i^*$, RMSD $r_{ik}$ and official
selected-pose PB validity $b_{ik}$,
\begin{equation}
 \mathrm{SR}=\frac{100}{M}\sum_i\mathbf1[r_{i,k_i^*}<2\,\text{\AA}],\quad
 \mathrm{SR}_{\rm PB}=\frac{100}{M}\sum_i\mathbf1[r_{i,k_i^*}<2\,\text{\AA}]b_{i,k_i^*}.
\end{equation}
RMSD uses symmetry-aware no-alignment heavy-atom CalcRMS. Solid bars show
$\mathrm{SR}_{\rm PB}$; hatched extensions show $\mathrm{SR}-\mathrm{SR}_{\rm PB}$.
Top-$k$ uses the full-bank confidence order, while generation-prefix coverage
uses saved generation indices. Both reach the RMSD-only oracle at 100; prefix
reselection is a distinct, potentially nonmonotone selected-pose endpoint.

The external cohorts are Astex Diverse Set (85), PoseBusters v2 (308),
PhiBench (206), FoldBench (558) and auxiliary OpenBind (925). PhiBench includes
three reconstructed cases; OpenBind includes two flagged noncovalent
approximations of covalent systems. FoldBench PB evaluation includes three
energy-reference InChI compatibility-repair shards. Original denominators are
retained. Report three-seed means and sample SD, except paired bootstrap
figures: average paired differences over seeds within each complex, then
resample complexes or exact-PDB groups 2,000 times and report percentile
95\% CIs in percentage points. Group resampling retains member complexes
and uses their complex-weighted mean; it is not protein-family clustering.

Sequence relatedness uses the 47,277 docking samples. For query binding
chains $Q$ and a training sample $h$, define
\begin{equation}
 s(Q,h)=100\frac{\max_{\pi}\sum_{q\in Q}n_{\rm identical}(q,\pi(q))}
 {\sum_{q\in Q}|q|},\qquad s_{\max}(Q)=\max_{h\in\mathcal H}s(Q,h),
\end{equation}
where $\pi$ is a one-to-one assignment to eligible binding chains within
that training sample. Unaligned residues, unmatched chains and unknown
residues do not contribute identical-residue evidence; all observed query
residues remain in the denominator. Ligand--protein overlap requires exact
canonical heavy-isomeric ligand identity and $s(Q,h)\ge70$ in the same
training sample, searching all exact-ligand witnesses. The threshold is
descriptive, not a certified pocket match or binary leakage definition.
Six-cohort composition adds Validation (1,076), without implying a matched
three-seed validation performance bank. Opened external results are
descriptive and do not select checkpoints or tune settings.
.
% Implementation references and full typed-energy parameter tables are in docs/methods.
\section{Methods}
\label{sec:effdock-methods}

\subsection{Task, data and information boundary}
EFF-Dock predicts ligand poses given a receptor, ligand chemistry and a
supplied pocket centre. The experiments are retrospective supplied-pocket
redocking: pocket centres may be reference-ligand-derived and receptors may
be holo structures. The reference pose supplies supervision, mapping and
evaluation labels; its coordinates do not enter the learned vector field,
confidence ranker or minimized refinement objective.

Training structures originate from PLINDER 2024-06/v2. A sample is a system
and ligand-instance pair. The preserved split contains 47,310 training and
1,076 validation IDs. The released docking fine-tune loader contains 47,277
samples; confidence uses 43,092 training and 1,035 validation samples. Their
training intersection is 43,067, with 4,210 docking-only and 25 confidence-only
samples. These independently filtered inventories are not nested and do not
represent the union of every predecessor checkpoint's training exposure.
Public membership records provide IDs, exclusions and SHA-256 hashes.

The historical external exclusion combined a ligand-match condition with a
pocket-community-match condition whose reference communities were identified
from benchmark PDB hits in the processed pool. It is distinct from the newer
strict split builder. No post-hoc exclusion changes the released benchmarks.

\subsection{Fragment representation and equivariant model}
Ligands are partitioned at eligible single, non-ring, non-planar-conjugated
rotatable bonds, retaining at least two heavy atoms per fragment after
singleton merging. Let $f(a)$ assign atom $a$ to fragment $f$, $T_f$ be its
centroid, $q_f$ a unit quaternion and $\ell_a$ its fixed local coordinate:
\begin{equation}
 T_f=\frac{1}{|f|}\sum_{a\in f}x_a,\qquad
 x_a(T,q)=R(q_{f(a)})\ell_a+T_{f(a)}.
 \label{eq:effdock-reconstruction}
\end{equation}
The heterogeneous graph contains ligand atoms, fragments, protein atoms and
residue virtual nodes. Static topology and dynamic protein--ligand contacts
within $5\,\text{\AA}$ condition six equivariant interaction layers, using
32 radial basis functions and spherical harmonics through degree two. Hidden
irreps are $384\times0e+32\times1o+32\times1e+16\times2e+16\times2o$,
with flattened width 736. Irrep notation alone does not establish full
reflection equivariance of stereochemical inputs; the pose contract uses
proper rotations and joint SE(3) transformations.

The model outputs atom vector fields $g_a$. Define $r_a=x_a-T_{f(a)}$.
Unit-weight Newton--Euler aggregation gives
\begin{align}
 v_f&=\frac{1}{|f|}\sum_{a\in f}g_a,&
 \tau_f&=\sum_{a\in f}r_a\times g_a,\\
 I_f&=\sum_{a\in f}(\|r_a\|^2I_3-r_ar_a^{\mathsf T}),&
 \omega_f&=I_f^+\tau_f.
 \label{eq:effdock-head}
\end{align}
The pseudoinverse retains inertia eigenvalues strictly above 1\% of the
largest eigenvalue, yielding an observable-rotation projector $P_f$.

\subsection{Conditional flow and training objective}
For supplied centre $c$, translation priors are
$T_f^0=c+\sigma\epsilon_f$, $\epsilon_f\sim\mathcal N(0,I_3)$;
initial orientations are uniform on SO(3). Conditional interpolation is
\begin{align}
 T_f^t&=(1-t)T_f^0+tT_f^1,&
 q_f^t&=\operatorname{SLERP}(q_f^0,q_f^1;t),\\
 v_f^*&=T_f^1-T_f^0,&
 \omega_f^*&=\operatorname{axisAngle}(q_f^1\otimes(q_f^0)^{-1}).
 \label{eq:effdock-flow-targets}
\end{align}
Quaternion updates use left multiplication. With $F$ fragments and
$F_\omega$ fragments having more than one atom, the flow terms are
\begin{equation}
 L_v=\frac{1}{3F}\sum_f\|v_f-v_f^*\|^2,\qquad
 L_\omega=\frac{1}{3F_\omega}\sum_{f:|f|>1}
 \|\omega_f-P_f\omega_f^*\|^2.
 \label{eq:effdock-flow-loss}
\end{equation}
A zero eligible count yields zero angular loss. The angular denominator is
not the number of observable eigenvectors. The released fragmenter avoids
singletons. Optional time weighting of these flow terms is disabled.

The atom auxiliary loss compares instantaneous velocities, not endpoint
coordinates. For $r_a^t=x_a^t-T_{f(a)}^t$,
\begin{equation}
 u_a=v_{f(a)}+\omega_{f(a)}\times r_a^t,\quad
 u_a^*=v_{f(a)}^*+\omega_{f(a)}^*\times r_a^t,\quad
 L_{\rm atom}=\frac{1}{3A}\sum_a\|u_a-u_a^*\|^2.
 \label{eq:effdock-atom-loss}
\end{equation}
The auxiliary target uses the unprojected angular field; an exactly
unobservable rotation induces zero atom motion.

A separate distance-geometry term uses the one-step endpoint
\begin{equation}
 \widehat T_f^1=T_f^t+(1-t)v_f,\qquad
 \widehat q_f^1=\operatorname{Exp}_q((1-t)\omega_f)\otimes q_f^t.
\end{equation}
Here $\operatorname{Exp}_q$ converts an axis-angle vector to a quaternion.
For unordered inter-fragment atom pairs $\mathcal P_b$ in complex $b$,
\begin{align}
 D_b&=\frac{1}{|\mathcal P_b|}\sum_{(a,c)\in\mathcal P_b}
 \left(\|\widehat x_a^1-\widehat x_c^1\|-\|x_a^1-x_c^1\|\right)^2,\\
 L_{\rm DG}&=\frac{\sum_{b:|\mathcal P_b|>0}t_b^2D_b}
 {\sum_{b:|\mathcal P_b|>0}t_b^2},\qquad
 L=L_v+8L_\omega+0.3L_{\rm atom}+3L_{\rm DG}.
 \label{eq:effdock-total-loss}
\end{align}
A zero denominator yields zero $L_{\rm DG}$. Distributed reductions reproduce
global fragment, atom and weighted-complex means for the respective terms;
these are not equal-weight-per-complex reductions throughout.

The released 50k fine-tune samples
\begin{align}
 p(t)&=0.80p_{\rm SF}(t)+0.10U(0,0.3)+0.10\delta_0,\\
 p_{\rm SF}&=0.02U(0,1)+0.98\operatorname{Law}
 [\operatorname{sigmoid}(0.8+1.7Z)],\quad Z\sim\mathcal N(0,1).
\end{align}
It uses fresh AdamW state, global batch 64 on four GPUs, peak learning rate
$2\times10^{-5}$, weight decay 0.01, gradient clipping 1, and EMA decay 0.999.
Warmup lasts 1,000 updates; cosine decay begins at 40,000 and reaches
$2\times10^{-6}$ at 50,000. The released docking state is terminal EMA.

\subsection{Sampling and confidence}
Main sampling uses $N=100$ candidates, $S=10$ steps, $\sigma=2\,\text{\AA}$
and a $10\,\text{\AA}$ crop. The late-power-three Euler grid and updates are
\begin{align}
 t_j&=1-(1-j/S)^3,\quad j=0,\ldots,S,\\
 T_f^{j+1}&=T_f^j+\Delta t_jv_{\theta,f}(t_j),\qquad
 q_f^{j+1}=\operatorname{Exp}_q(\Delta t_j\omega_{\theta,f}(t_j))\otimes q_f^j.
 \label{eq:effdock-sampler}
\end{align}
Raw and refined banks preserve generation indices and are independently scored.

The confidence scorer uses four equivariant layers, paired docking features,
and attention pooling over global and contact features. It predicts log-one-plus pose
RMSD $\widehat y_k$, pose-success logit $z_k$, atom log displacement
$\widehat y_{ka}$ and atom-success logit $z_{ka}$. Atom and pose regression
use Huber loss with transition one; classification uses BCE with logits and
strict $2\,\text{\AA}$ labels. Let $y_k=\log(1+r_k)$:
\begin{align}
 L_{\rm rank}&=\underset{y_j-y_i>0.3}{\operatorname{mean}}
 \max(0,0.05-(\widehat y_j-\widehat y_i)),\\
 q_k&=\operatorname{sigmoid}((2-r_k)/0.2),\quad
 p_k=\frac{q_k}{\sum_jq_j},\qquad
 L_{\rm list}=-\sum_kp_k\log\operatorname{softmax}(z)_k,\\
 L_{\rm conf}&=0.2L_{\rm atom}+0.2L_{\rm atom\mbox{-}ok}
 +0.3L_{\rm pose}+0.4L_{\rm pose\mbox{-}ok}
 +0.1L_{\rm rank}+L_{\rm list}.
 \label{eq:effdock-confidence}
\end{align}
Here the confidence atom term is displacement regression, distinct from the
docking atom-velocity term in Eq.~\ref{eq:effdock-atom-loss}. Empty eligible
pairs give zero rank loss. Numerically zero quality sums use a uniform target;
one-candidate pools have zero listwise loss.

Training banks contain 32 raw poses, 32 refined poses and a crystal anchor;
the loader samples a bounded subset. Validation excludes crystal anchors.
The continuation uses Muon at $2\times10^{-4}$ plus AdamW at $3\times10^{-6}$,
effective global batch two, seed 45, 2,000-update warmup and cosine decay over
the final 50,000 of 100,000 updates, ending at 0.05 of the peak rates.
U70k was selected only on 1,035 internal complexes, scoring 622/1,035 (60.10\%)
Top-1 RMSD success, compared with 617/1,035 at U100k.

\subsection{Energy-based rigid refinement}
Refinement is a separate post-sampling operation with a fixed receptor:
\begin{equation}
 E=E_{\rm physical}+E_{\rm interaction},\qquad F_a=-\nabla_{x_a}E.
 \label{eq:effdock-energy}
\end{equation}
The physical terms are harmonic bond/angle restraints, periodic proper
torsions, planar/chiral impropers, softened Lennard--Jones interactions,
compact protein--ligand overlap barriers and repulsive receptor obstacles.
Physical parameters use EFF-FF version 2.2.0 and interaction parameters version
1.6.0. The chemical-constraint channel has scale zero, while the physical
chiral-improper scale is one. The input conformer's geometry, not the target
crystal pose, defines restraint reference values.
Bond restraints act on cut bonds, angles and impropers span fragments, and
proper-torsion weights sum to one per cut bond over its retained quadruplets.
\begin{align}
 E_{\rm bond}&=\tfrac12\sum_bk_b(d_b-d_b^0)^2,&
 E_{\rm angle}&=\tfrac12\sum_ak_a(\theta_a-\theta_a^0)^2,\\
 E_{\rm proper}&=\sum_pk_pw_p[1+\cos(n_p\phi_p-\delta_p)],\\
 E_{\rm improper}&=\tfrac12\sum_{i\in\rm planar}k_i[1-\cos(2\phi_i)]
 +\tfrac12\sum_{i\in\rm chiral}k_i\operatorname{wrap}(\phi_i-\phi_i^0)^2.
\end{align}
For $Q(u)=u^3(10-15u+6u^2)$, $S(d)=Q(\operatorname{clip}((8-d)/2,0,1))$
and $\rho=\sqrt{d^2+0.75^2}$,
\begin{equation}
 E_{{\rm LJ},ij}=s_{ij}S(d_{ij})D_{ij}
 [(X_{ij}/\rho_{ij})^{12}-2(X_{ij}/\rho_{ij})^6].
\end{equation}
Pair atom parameters use geometric mixing; the versioned table specifies
exclusions and the default 1--4 scale of 0.5. The steric term is
\begin{align}
 d_{ij}^{\rm safe}&=0.8(r_i^{\rm vdW}+r_j^{\rm vdW}),&
 p_{ij}&=0.1\log[1+e^{(d_{ij}^{\rm safe}-d_{ij})/0.1}],\\
 E_{\rm steric}&=10\sum_{ij}p_{ij}^2
 Q\!\left(\operatorname{clip}\left(\frac{d_{ij}^{\rm safe}+0.5-d_{ij}}{0.5},0,1\right)\right).
\end{align}
The interaction terms are hydrophobic, hydrogen bond, screened formal charge,
pi stacking, cation--pi, halogen bond and profile-specific metal coordination.
For bounded pair weights $w_{ij}$, site saturation is
\begin{equation}
 O_i=1-\prod_j[1-(1-10^{-7})w_{ij}],\qquad E_{\rm site}=-\epsilon\sum_iO_i.
\end{equation}
Symmetric ring-system saturation averages occupancies in both directions.
Screened formal-charge energy uses
\begin{equation}
 E_{\rm charge}=\sum_{ij}\frac{332.06371}{78.5}q_iq_j
 \frac{e^{-0.127\rho_{ij}}}{\rho_{ij}}S_q(d_{ij}^2),\quad
 \rho_{ij}=\sqrt{d_{ij}^2+0.75^2},
\end{equation}
with the quintic switch between squared distances $8^2$ and $10^2$.
The parameter tables and implementation specify all typed geometric gates,
metal profiles and repulsive-only obstacle terms. These diagnostic energies
are not calibrated binding free energies.

Energy-force projection uses atomic masses, unlike the learned head. For
fragment mass $M_f$, centre of mass $C_f$, centroid $T_f$ and $r_a=x_a-C_f$,
\begin{align}
 C_f&=\frac{\sum_{a\in f}m_ax_a}{M_f},&
 I_f&=\sum_{a\in f}m_a(\|r_a\|^2I_3-r_ar_a^{\mathsf T}),\\
 \Omega_f&=I_f^+\sum_{a\in f}r_a\times F_a,&
 V_f&=\frac{\sum_{a\in f}F_a}{M_f}+\Omega_f\times(T_f-C_f).
 \label{eq:effdock-energy-projection}
\end{align}
Independent per-pose backtracking tries $\alpha=0.5^b$, $b=0,\ldots,12$,
accepting finite energy no greater than
$E+10^{-10}\max(1,|E|)$. Each step is capped at $0.10\,\text{\AA}$
translation, five degrees rotation and $0.10\,\text{\AA}$ actual atom
movement. At most 100 iterations are run. Stop after 20 consecutive accepted
steps with maximum atom movement below $10^{-5}\,\text{\AA}$, or, from
iteration 25, five consecutive steps with energy decrease at most
$0.02+0.001\max(1,|E|)$. External runs use an $18\,\text{\AA}$ receptor
shell, $8\,\text{\AA}$ physical cutoff and $10\,\text{\AA}$ maximum
interaction cutoff. Training-bank refinement instead used displacement
threshold $0.01\,\text{\AA}$ with patience five. Failed line searches retain
the last accepted pose and an explicit status.

\subsection{Guided-sampling diagnostics}
The separately captioned guidance/budget and pocket/prior figures use
direct ODE guidance with $\eta=2$; the main condition has $\eta=0$.
Let $U(v,\omega)$ be the induced atom velocity and
$\operatorname{RMS}(U)=\sqrt{A^{-1}\sum_a\|U_a\|^2}$. Before capping,
\begin{equation}
 (\Delta v,\Delta\omega)=\eta\bar r_j
 \frac{\operatorname{RMS}(U(v_\theta,\omega_\theta))}
 {\operatorname{RMS}(U(V,\Omega))}(V,\Omega),\qquad
 \bar r_j=\frac{1}{\Delta t_j}\int_{t_j}^{t_{j+1}}\max(0,2t-1)\,dt.
\end{equation}
Here $(V,\Omega)$ projects energy forces capped at atom norm 20. The physical
LJ softcore decreases from 1.5 to $0.75\,\text{\AA}$ over the active
half-interval, evaluated at the active interval midpoint. The drift vanishes
when either RMS is at most $10^{-8}$ or the direction is nonfinite.
One scale per pose enforces added translation/angular velocity limits of
5 and 5 per unit flow time, and an estimated atom-displacement bound of
$0.25\,\text{\AA}$. The sampler adds the drift to the learned field before
its Euler update. Chemical-constraint drift is zero. This is separate from
the subsequent energy-based refinement.

\subsection{Selection, metrics and relatedness}
Predicted RMSD is $\widehat r_k=\max(0,\operatorname{expm1}
(\operatorname{clip}(\widehat y_k,-2,5)))$. Ordinary selection minimizes
$(\widehat r_k,k)$ lexicographically. Chirality filtering restricts this
minimum to candidates consistent with input tetrahedral stereochemistry,
falling back to ordinary selection if none passes. The main mask contains no
reference RMSD, PB test or E/Z criterion.

For $M$ complexes, selected index $k_i^*$, RMSD $r_{ik}$ and official
selected-pose PB validity $b_{ik}$,
\begin{equation}
 \mathrm{SR}=\frac{100}{M}\sum_i\mathbf1[r_{i,k_i^*}<2\,\text{\AA}],\quad
 \mathrm{SR}_{\rm PB}=\frac{100}{M}\sum_i\mathbf1[r_{i,k_i^*}<2\,\text{\AA}]b_{i,k_i^*}.
\end{equation}
RMSD uses symmetry-aware no-alignment heavy-atom CalcRMS. Solid bars show
$\mathrm{SR}_{\rm PB}$; hatched extensions show $\mathrm{SR}-\mathrm{SR}_{\rm PB}$.
Top-$k$ uses the full-bank confidence order, while generation-prefix coverage
uses saved generation indices. Both reach the RMSD-only oracle at 100; prefix
reselection is a distinct, potentially nonmonotone selected-pose endpoint.

The external cohorts are Astex Diverse Set (85), PoseBusters v2 (308),
PhiBench (206), FoldBench (558) and auxiliary OpenBind (925). PhiBench includes
three reconstructed cases; OpenBind includes two flagged noncovalent
approximations of covalent systems. FoldBench PB evaluation includes three
energy-reference InChI compatibility-repair shards. Original denominators are
retained. Report three-seed means and sample SD, except paired bootstrap
figures: average paired differences over seeds within each complex, then
resample complexes or exact-PDB groups 2,000 times and report percentile
95\% CIs in percentage points. Group resampling retains member complexes
and uses their complex-weighted mean; it is not protein-family clustering.

Sequence relatedness uses the 47,277 docking samples. For query binding
chains $Q$ and a training sample $h$, define
\begin{equation}
 s(Q,h)=100\frac{\max_{\pi}\sum_{q\in Q}n_{\rm identical}(q,\pi(q))}
 {\sum_{q\in Q}|q|},\qquad s_{\max}(Q)=\max_{h\in\mathcal H}s(Q,h),
\end{equation}
where $\pi$ is a one-to-one assignment to eligible binding chains within
that training sample. Unaligned residues, unmatched chains and unknown
residues do not contribute identical-residue evidence; all observed query
residues remain in the denominator. Ligand--protein overlap requires exact
canonical heavy-isomeric ligand identity and $s(Q,h)\ge70$ in the same
training sample, searching all exact-ligand witnesses. The threshold is
descriptive, not a certified pocket match or binary leakage definition.
Six-cohort composition adds Validation (1,076), without implying a matched
three-seed validation performance bank. Opened external results are
descriptive and do not select checkpoints or tune settings.
